\renewcommand{\abstractname}{Introduction}

\begin{abstract}
    Le présent rapport s'inscrit dans le cadre d'une alternance d'un an en contrat de professionnalisation, effectuée par un étudiant de l'INSA Centre Val-de-Loire au sein du \gls{cea}, sur le site de l'\gls{ines} pour le \gls{lela}.

    L'alternance vise à consolider d'une part les connaissances du laboratoire en matière de méthodes d'identification des paramètres de modèles de bâtiment, et d'autre part à poursuivre le développement de modèles d'optimisation de ces paramètres selon plusieurs paradigmes ({\gls{algorithm}e}s méta-heuristiques, apprentissage automatique, etc.).

    Ce livrable propose de se focaliser sur les éléments non scientifiques (dits aspects ``humanités'') de l'alternance et sera complété par un autre rapport (dit ``rapport technique'') ce semestre puis par un rapport global à la fin de la période d'alternance.

    La première partie présentera le CEA à trois niveaux hiérarchiques différents, où seront rappelés pour chacun d'entre eux quelques éléments historiques, les missions confiées et l'organisation des équipes.
    Fort de ces informations, nous conclurons chaque présentation par une analyse sommaire de la stratégie de recherche et d'innovation mise en place par le CEA…

    La seconde partie synthétisera les premières semaines d'activité de l'alternant, en détaillant le parcours d'intégration au sein des différentes unités du CEA, les missions confiées \textit{(terminées, en cours et à venir)}, ainsi que les compétences apportées par l'alternant et celles développées durant cette période.
\end{abstract}

